\documentclass[11pt]{amsart}
\usepackage{amsmath, amssymb, amsthm}
\usepackage[margin=1in]{geometry}
\usepackage[hypertexnames=false]{hyperref}
\hypersetup{hidelinks}
\usepackage{booktabs}
\usepackage{enumitem}

\theoremstyle{plain}
\newtheorem{theorem}{Theorem}
\newtheorem{corollary}[theorem]{Corollary}
\newtheorem{proposition}[theorem]{Proposition}
\newtheorem{lemma}[theorem]{Lemma}
\newtheorem*{conjecture-fwd}{Conjecture (strict form)}
\newtheorem*{hypothesis}{Hypothesis H3}

\theoremstyle{definition}
\newtheorem{conjecture}[theorem]{Conjecture}

\newcommand{\RESULTabstract}{the only unitary perfect numbers with $a\le 17$ or $a\in\{19, 21, 23, 25\}$ are $6$, $60$, $90$ and $87360$; so a sixth unitary perfect number has $a\ge 18$ and $a\notin\{19, 21, 23, 25\}$.}
\newcommand{\RESULTtheorem}{Let $n=2^am$ be a unitary perfect number with $m$ odd.  If $a\le 17$ or $a\in\{19, 21, 23, 25\}$, then $n$ is one of $6$, $60$, $90$ and $87360$.  Consequently a sixth unitary perfect number has $a\ge 18$ and $a\notin\{19, 21, 23, 25\}$.}
\newcommand{\RESULTopen}{For example, at $a=22$ and $a=29$ the search reaches branches in which the next component would have to be sought among all primes up to about $10^{13}$.}
\newcommand{\Aall}{17}
\newcommand{\Apy}{21}
\newcommand{\RESULTwallproof}{}
\newcommand{\RESULTwallabstract}{The search recovers the four known unitary perfect numbers with $a\le 6$, and with the component $5^4$ fixed it finds Wall's $24$-digit example at $a=18$.}
\newcommand{\RESULTwallcheck}{The search recovers the four known unitary perfect numbers with $a\le6$.  Run at $a=18$ with the component $5^4$ placed first, it finds Wall's $11$-component example in $488$ nodes.}

\title[Unitary perfect numbers with a small power of two]{Unitary perfect numbers with a small power of two, and a conditional bounded-box reduction}
\author[T. Maciejewski]{Tom Maciejewski}
\address{Independent researcher, New Jersey, USA}
\email{tommaciejewski4@gmail.com}
\thanks{ORCID: 0009-0004-5702-5302.}
\date{November 2026}

\subjclass[2020]{Primary 11A25; Secondary 11Y05, 11N69, 11Y50, 11A41}
\keywords{Unitary perfect numbers; Subbarao--Warren conjecture;
exhaustive search; 3-Higgs primes; cyclotomic primitive divisors;
Cunningham numbers}

\begin{document}

\begin{abstract}
A unitary perfect number is a positive integer $n$ with $\sigma^*(n)=2n$,
where $\sigma^*(n)$ is the sum of the divisors $d$ of $n$ with
$\gcd(d,n/d)=1$.  Only five are known, and it has been conjectured that there
are only finitely many.  Every unitary perfect number has the form $n=2^am$ with $m$ odd
and $a\ge1$.  We prove that \RESULTabstract{}  The previously published bound
was $a\ge 11$.

The proof is an exhaustive search over the odd prime-power components of
$m$, taken in increasing order, which uses nothing but the balance
$(2^a+1)\prod(p^e+1)=2^{a+1}\prod p^e$: every odd prime must occur to the same
power on both sides, and the power of $2$ allows at most $a+1$ components.
\RESULTwallabstract{}

This version corrects the earlier ones.  Version~1 asserted, and
Proposition~1 of version~2 claimed to prove, that every prime divisor of a
unitary perfect number is a 3-Higgs prime.  That proof is invalid and the
statement is open, so the bounded-box elimination of ``impostor'' kernels for
$a\le 10000$ claimed there holds only under it.  We state the hypothesis
explicitly, repair a second gap (a 2-adic budget step assumed that
$v_2(p^e+1)$ cannot decrease as $e$ grows), settle unconditionally every case
in which that step had been used, and withdraw several conjectures and claims
of version~2 about the numbers $2^m+1$ that were false or unsupported.
\end{abstract}

\maketitle

\section{Introduction}

For a positive integer $n$ write $\sigma^*(n)=\sum_{d\,\Vert\,n}d$, where
$d\,\Vert\,n$ means $d\mid n$ and $\gcd(d,n/d)=1$.  A \emph{unitary perfect
number} is an $n$ with $\sigma^*(n)=2n$.  The known ones are
\[
6,\quad 60,\quad 90,\quad 87360,\quad 146361946186458562560000 .
\]
Subbarao and Warren~\cite{SW1966} introduced them and found the first four;
Wall~\cite{Wall1975} found the fifth.

\begin{conjecture-fwd}
The five integers above are the only unitary perfect numbers.
\end{conjecture-fwd}

Every unitary perfect number is even~\cite{SW1966}; write $n=2^am$ with
$m$ odd.  What is known about a sixth one:
\begin{itemize}[leftmargin=*]
\item $a\ge 11$ (Subbarao~\cite{Subbarao1970}; see also~\cite{SCNW1972} and
\cite{Wall1987,Wall1988});
\item $m$ has at least nine distinct prime factors (Wall~\cite{Wall1988});
\item some prime power $p^e\,\Vert\,m$ exceeds $2^{15}$ (Wall~\cite{Wall1987});
\item $m$ is not squarefree (Graham~\cite{Graham1989}).
\end{itemize}

\subsection*{Main result}
Our main theorem extends the first item.

\begin{theorem}\label{thm:intro}
\RESULTtheorem{}
\end{theorem}

Theorem~\ref{thm:intro} is proved in Section~\ref{sec:search} by an exhaustive
search, one value of $a$ at a time, that uses only the balance
\begin{equation}\label{eq:ast}
(2^a+1)\prod_{p^e\,\Vert\, m}(p^e+1)=2^{a+1}\prod_{p^e\,\Vert\, m}p^e .
\end{equation}
It keeps exact track of how often each odd prime has been produced on the left
and used on the right, and of how much of the power $2^{a+1}$ has been spent.
None of the published bounds listed above is used, and the search finds the
known examples, so it also serves as an independent check on them.

\subsection*{Correction to the earlier versions}
Versions~1 and~2 of this paper (arXiv:2605.20475v1 and v2, titled
\emph{Bounded-box reductions in the Subbarao--Warren problem for unitary perfect
numbers}) stated that every prime divisor of a unitary perfect number is a
3-Higgs prime (defined in Section~\ref{sec:higgs}), and built on that
statement.  Version~1 gave no proof.  Proposition~1 of version~2 gave one, but
it used the false divisibility $p-1\mid p^e+1$ (in fact
$p^e+1\equiv 2\pmod{p-1}$), and we withdraw it.  We know no proof: the
statement holds for the five known examples, but nothing in~\eqref{eq:ast}
appears to force it.  We now call it \emph{Hypothesis~H3}.  Three
consequences:
\begin{enumerate}[label=(\arabic*),leftmargin=*]
\item The elimination of the five impostor kernels for $1\le a\le10000$
(Theorem~3 and Corollary~4 of version~2) holds under Hypothesis~H3 and is
restated as Theorem~\ref{thm:conditional}.  The classification of source
kernels that produced those five kernels was not specified precisely enough to
be reproduced, and we no longer rely on it.
\item The 2-adic overshoot filter of version~2 was invalid even under H3: its
monotonicity lemma assumed that raising an exponent $e$ cannot lower
$v_2(p^e+1)$, but $v_2(p^e+1)=1$ whenever $e$ is even.  That filter was used for
ten candidates.  All ten are now settled unconditionally
(Proposition~\ref{prop:ten}).
\item Section~\ref{sec:Heven} keeps version~2's unconditional results on
$H_{\mathrm{even}}$, which are statements about the numbers $2^m+1$, and drops
its conjectural programme, which contained false statements
(Section~\ref{sec:changes}).  Their bearing on unitary perfect numbers goes
through Hypothesis~H3 only.
\end{enumerate}
A list of changes is in Section~\ref{sec:changes}.

\section{Setup}\label{sec:setup}

Let $n=2^am$ be unitary perfect, with $m$ odd and $a\ge1$.  An \emph{odd
component} of $n$ is a prime power $p^e\,\Vert\, m$.  Since $\sigma^*$ is
multiplicative with $\sigma^*(p^e)=p^e+1$, the equation $\sigma^*(n)=2n$ is the
balance~\eqref{eq:ast}.  Comparing valuations on its two sides gives, for every
odd prime $q$,
\begin{equation}\label{eq:exponent}
v_q(m)=v_q(2^a+1)+\sum_{p^e\,\Vert\, m}v_q(p^e+1),
\end{equation}
and, for the prime $2$,
\begin{equation}\label{eq:budget}
a+1=\sum_{p^e\,\Vert\, m}v_2(p^e+1).
\end{equation}
Each term of~\eqref{eq:budget} is at least $1$, so $m$ has at most $a+1$ odd
components.  We call $2^a+1$ the \emph{seed}: by~\eqref{eq:exponent} each of its
prime factors divides $m$.  (Versions~1 and~2 wrote~\eqref{eq:exponent} with
$v_q(p+1)$ in place of $v_q(p^e+1)$ and without the seed term.)

\section{An exact search over the power of two}\label{sec:search}

Throughout this section $n=2^a m$ is a unitary perfect number with $m$ odd and
$a\ge 1$ (every unitary perfect number is even~\cite{SW1966}).  Write
$m=c_1c_2\cdots c_k$, where the \emph{components} $c_i=p_i^{e_i}$ are powers of
distinct odd primes, ordered so that $c_1<c_2<\cdots<c_k$.  Then
$\sigma^*(n)=(2^a+1)\prod_i(c_i+1)$, and $\sigma^*(n)=2n$ is the balance
\begin{equation}\label{eq:balance}
(2^a+1)\prod_{i=1}^{k}(c_i+1)=2^{a+1}\prod_{i=1}^{k}c_i .
\end{equation}

\begin{lemma}\label{lem:balance}
\begin{enumerate}[label=(\roman*),leftmargin=*]
\item $\prod_{i=1}^{k}\bigl(1+1/c_i\bigr)=T_a:=2^{a+1}/(2^a+1)$.
\item For every odd prime $q$,
$v_q(2^a+1)+\sum_{i=1}^{k}v_q(c_i+1)=v_q(m)$; in particular the left side
is $e_i$ if $q=p_i$ and $0$ if $q\nmid m$.
\item $\sum_{i=1}^{k}v_2(c_i+1)=a+1$. Since $v_2(c+1)\ge 1$ for odd $c$,
$k\le a+1$.
\end{enumerate}
\end{lemma}

\begin{proof}
Divide \eqref{eq:balance} by $(2^a+1)\prod c_i$ for (i), and compare
$q$-adic and $2$-adic valuations of its two sides for (ii) and (iii), using that
$2^a+1$ is odd.
\end{proof}

Part (ii) corrects the exponent formula of the earlier versions, which
omitted the seed term $v_q(2^a+1)$.

\subsection{Prefix invariants}
For $0\le j\le k$ put
\[
F_j=\prod_{i\le j}\Bigl(1+\frac1{c_i}\Bigr),\quad
S_j=\sum_{i\le j}v_2(c_i+1),\quad
R_j=a+1-S_j,\quad
N_j(q)=v_q(2^a+1)+\sum_{i\le j}v_q(c_i+1),
\]
let $D_j=\{p_1,\dots,p_j\}$, and let
$Q_j=\{q\text{ odd prime}: N_j(q)>0,\ q\notin D_j\}$ be the primes that are
\emph{demanded} but not yet placed.

\begin{lemma}\label{lem:prefix}
For every $0\le j<k$:
\begin{enumerate}[label=(\alph*),leftmargin=*]
\item $F_j<T_a$, and $F_j(1+1/c_{j+1})\le T_a$.
\item $N_j(p_i)\le e_i$ for every $i\le j$.
\item $Q_j\subseteq\{p_{j+1},\dots,p_k\}$; hence $|Q_j|\le k-j\le R_j$.
\item If $p_{j+1}\in Q_j$ then $e_{j+1}\ge N_j(p_{j+1})$.  If
$p_{j+1}\notin Q_j$ then $|Q_j|+1\le R_j$.
\item If $p_{j+1}\notin Q_j$ and $|Q_j|=R_j-1$, then $v_2(c_{j+1}+1)=1$ and
every odd prime factor of $c_{j+1}+1$ lies in $D_j\cup Q_j$.
\item Let $m_q$ be the least power $q^e>c_j$ with $e\ge N_j(q)$, for
$q\in Q_j$, and let $f_j=R_j-|Q_j|$. Then
\[
\frac{T_a}{F_j}\le\Bigl(1+\frac1{c_{j+1}}\Bigr)
\prod_{q\in Q_j}\Bigl(1+\frac1{\max(c_{j+1}+1,\,m_q)}\Bigr)
\Bigl(1+\frac1{c_{j+1}+1}\Bigr)^{f_j}.
\]
\end{enumerate}
\end{lemma}

\begin{proof}
(a) $T_a/F_j=\prod_{i>j}(1+1/c_i)$ by Lemma~\ref{lem:balance}(i), and each
factor exceeds $1$.  (b) $N_j(p_i)\le N_k(p_i)=e_i$ by
Lemma~\ref{lem:balance}(ii).  (c) A prime with $N_j(q)>0$ divides $m$ by
Lemma~\ref{lem:balance}(ii), and is not among $p_1,\dots,p_j$ when $q\in Q_j$.
Each later component costs $v_2(c_i+1)\ge1$, and these costs add up to $R_j$
by Lemma~\ref{lem:balance}(iii).  (d) The first claim is (b) at $j=k$. For the
second, the later components then include one for each prime of $Q_j$ and one
for $p_{j+1}$.  (e) After step $j+1$ the demanded set contains $Q_j$ and every
odd prime factor of $c_{j+1}+1$ outside $D_j\cup Q_j$ (none of them is
$p_{j+1}$), and by (c) its size is at most
$R_{j+1}=R_j-v_2(c_{j+1}+1)$.  So $R_j-1+t\le R_j-v_2(c_{j+1}+1)$, where $t$ is
the number of such primes, which forces $v_2(c_{j+1}+1)=1$ and $t=0$.
(f) The later components are $c_{j+1}$, one power of each $q\in Q_j$ other
than $p_{j+1}$, which exceeds both $c_{j+1}$ and $c_j$ and so is at least
$\max(c_{j+1}+1,m_q)$, and at most $f_j$ further components, each larger than
$c_{j+1}$.  The right side bounds their product, with one factor to spare.
\end{proof}

\subsection{The search}
Fix $a$.  The search builds the sequence $c_1<c_2<\cdots$ one component at
a time.  At a prefix of length $j$ it computes $F_j$, $S_j$, $N_j$ and $Q_j$
exactly, and takes as candidates for $c_{j+1}$ the prime powers $p^e$ with
\[
c_j<p^e,\qquad F_j(1+p^{-e})\le T_a,\qquad
\text{the inequality of Lemma~\ref{lem:prefix}(f) holds},
\]
subject to (b)--(e).  Concretely the candidates are: the powers $q^e$ with
$q\in Q_j$ and $e\ge N_j(q)$; if $f_j=1$, the numbers $c=2M-1$ with $M$ odd and
composed of primes in $D_j\cup Q_j$ (within the capacities left by (b)) that are
powers of a new prime (Lemma~\ref{lem:prefix}(e)); and if $f_j\ge2$, all powers
of new primes in the range.  A candidate is kept only if, after it is added,
(b) and (c) still hold.  The prefix is a solution when $F_j=T_a$,
$S_j=a+1$ and $N_j(q)$ equals the exponent of $q$ for every prime $q$.

\begin{theorem}\label{thm:complete}
For each $a$, every unitary perfect number $2^a m$ with $m$ odd is found by the
search, provided no branch is reported as open.
\end{theorem}

\begin{proof}
Let $2^a m$ be unitary perfect with components $c_1<\cdots<c_k$.  By
Lemma~\ref{lem:prefix}, for each $j<k$ the component $c_{j+1}$ is among the
candidates at the prefix $(c_1,\dots,c_j)$ and survives every test applied to
it, and the full sequence satisfies the solution conditions by
Lemma~\ref{lem:balance}.  The search leaves a candidate unexamined or
undecided in only three situations, and each is reported as an open branch:
when $f_j\ge2$ and the range for $c_{j+1}$ extends beyond the limit up to
which new primes are scanned ($2\cdot10^{10}$, and $3\cdot10^{10}$ for
$a=16$); when $c+1$ cannot be factored completely and the tests that need
only part of its factorization do not exclude $c$; and, in the C program,
when a candidate is a power of a prime above $2^{62}$.
\end{proof}

\subsection{Results}

\begin{proof}[Proof of Theorem~\ref{thm:intro}]
For each $a$ in the statement, the search of Theorem~\ref{thm:complete}
terminates with no open branch.  It returns exactly the components of $6$ and
$90$ for $a=1$, of $60$ for $a=2$, of $87360$ for $a=6$\RESULTwallproof{}, and
nothing for the other values.  Table~\ref{tab:search} gives the size of each
search.
\end{proof}

\begin{table}[h]
\centering
\small
\begin{tabular}{rrr@{\qquad}rrr@{\qquad}rrr}
\toprule
$a$ & nodes & sol. & $a$ & nodes & sol. & $a$ & nodes & sol.\\
\midrule
1 & 2 & 2 & 8 & 32 & 0 & 15 & 146 & 0\\
2 & 3 & 1 & 9 & 12 & 0 & 16 & 713\,484\,870 & 0\\
3 & 2 & 0 & 10 & 28 & 0 & 17 & 606 & 0\\
4 & 5 & 0 & 11 & 54 & 0 & 19 & 1\,811 & 0\\
5 & 7 & 0 & 12 & 557 & 0 & 21 & 7\,016 & 0\\
6 & 12 & 1 & 13 & 121 & 0 & 23 & 8\,951\,378 & 0\\
7 & 13 & 0 & 14 & 5\,897 & 0 & 25 & 440\,422\,296 & 0\\
\bottomrule
\end{tabular}
\medskip
\caption{The exact search for each settled $a$: number of nodes visited and number of solutions found (no branch is open in any of these runs). The solutions are $6$ and $90$ ($a=1$), $60$ ($a=2$) and $87360$ ($a=6$).}
\label{tab:search}
\end{table}

The previously published bound is $a\ge11$, due to Subbarao, Cook, Newberry
and Weber~\cite{SCNW1972}, as reported by Wall~\cite{Wall1987,Wall1988}.  Wall's
results that a sixth unitary perfect number has at least nine odd
components~\cite{Wall1988} and an odd component exceeding
$2^{15}$~\cite{Wall1987} are not used.

\subsection*{Implementation and checks}
The search was implemented twice: in Python with exact rational arithmetic,
and as a C port using GMP.  The two programs visit exactly the same number of
nodes on every input where both were run: all $a\le15$, $a=17,19,21$, and the
kernel runs of Proposition~\ref{prop:ten} at $a=45$ and $63$ and with $5^4$
placed first at $a=18$.  Factorizations of numbers up to $2\cdot10^7+2$ are read
from a table of smallest prime factors.  New primes between $2\cdot10^7$ and
$2\cdot10^{10}$ are found with a segmented sieve; for each, $c+1$ is divided by
the primes already present and the cofactor tested, exactly as in
Lemma~\ref{lem:prefix}(b),(c).  This path was checked by rebuilding the C program
with the table limited to $10^5$: it visits the same $8{,}951{,}378$ nodes at
$a=23$ as the standard build.  Other numbers are factored by trial division
and Pollard's $\rho$, and a branch whose numerator cannot be factored
completely is reported as open.  An open branch can be resumed from its node
with a larger scan limit, since the node determines the components already
chosen and the lower bound for the next one.  For $a=16$ the scan limit was
set to $3\cdot10^{10}$ from the start, so that no branch is left open.  Every
pruning decision is made in exact integer arithmetic: a floating-point estimate
only chooses where the binary search for the end of a candidate range starts,
and that end is then certified exactly.  In the C program every prime that is used is below $2^{62}$, where the
Baillie--PSW test is known to be correct; otherwise a probable-prime test is used
only to reject candidates, and a ``composite'' verdict is always correct.  The
only larger prime in any run is the $22$-digit factor
$5302306226370307681801$ of $2^{117}+1$ (Proposition~\ref{prop:ten}), whose
primality was proved with PARI/GP.  \RESULTwallcheck{}

\section{The 3-Higgs hypothesis and the bounded box}\label{sec:higgs}

This section restates the main result of the earlier versions with its
dependence on an unproven hypothesis made explicit, and repairs one of its
steps.

\subsection{Hypothesis H3}
A prime $p$ is \emph{3-Higgs} (OEIS A057447~\cite{OEIS}; see also
Burris--Yeats~\cite{BurrisYeats}) if $p-1$ divides the cube of the product of
the smaller 3-Higgs primes; equivalently, every prime $q\mid p-1$ is 3-Higgs and
$v_q(p-1)\le 3$.  Let $\mathcal P_3$ be the set of 3-Higgs primes.  The smallest
odd prime not in $\mathcal P_3$ is $17$.

\begin{hypothesis}
Every prime divisor of a unitary perfect number is 3-Higgs.
\end{hypothesis}

Every prime divisor of the five known unitary perfect numbers is 3-Higgs.  We
know of no proof of Hypothesis~H3 and no argument that makes it plausible beyond
that evidence; it is used only in this section and in the parts of
Section~\ref{sec:Heven} that say so.

\subsection{Kernels}
For a unitary perfect number $n$, let $G_n$ be the directed graph on the odd
primes dividing $n$, with an edge $p\to r$ whenever $r\mid p^{v_p(n)}+1$.  A
\emph{source kernel} is a strongly connected set $K$ of at least two such primes
with no edge entering $K$ from outside; for $q\in K$, \eqref{eq:exponent} then
gives $v_q(n)=v_q(2^a+1)+\sum_{p\in K}v_q(p^{v_p(n)}+1)$.  Version~2 reported
that a search over such kernels in the box
\[
\mathcal B := \{ p \le 2000,\ e \le 6,\ p^e \le 10^9,\ |K| \le 6 \},
\]
using the filters below, leaves, besides kernels of the known examples, the
following five \emph{impostor} kernels.  That search was not specified
precisely enough to be reproduced, and we use the list only as a list of cases.
\begin{center}\begin{tabular}{lll}
\toprule
Kernel & Forced exponents & Seed congruence \\
\midrule
$3^2\,5^3$ & $\{3{:}2,\,5{:}3\}$ & $a \equiv 10 \pmod{20}$ \\
$3^4\,41$ & $\{3{:}4,\,41{:}1\}$ & $a \equiv 9 \pmod{18}$ \\
$5^2\,13^2$ & $\{5{:}2,\,13{:}2\}$ & $a \equiv 6 \pmod{12}$ \\
$5^4\,157^2\,313$ & $\{5{:}4,\,157{:}2,\,313{:}1\}$ & $a \equiv 130 \pmod{260}$ \\
$5^4\,29\,157^2\,313$ & $\{5{:}4,\,29{:}1,\,157{:}2,\,313{:}1\}$ & $a \equiv 26 \pmod{52}$ \\
\bottomrule
\end{tabular}\end{center}

\subsection{The filters}
For an impostor kernel $K$ and an $a$ in its seed class, the \emph{candidate}
$(K,a)$ asks for a unitary perfect number with $2^a\,\Vert\, n$ in which every
prime power of $K$ is an odd component.

\emph{Filter Z (Zsigmondy).}  For $e>1$, the number $p^e+1$ has a prime
divisor $r$ that divides no $p^j+1$ with $j<e$ (Zsigmondy~\cite{Zsigmondy}; the
only exception, $2^3+1$, does not arise here), and such an $r$ satisfies
$r\equiv1\pmod{2e}$.  Under H3, $r$ is 3-Higgs,
so every prime factor of $2e$ is 3-Higgs and occurs in $2e$ at most three times.
Applied to the seed exponent, this excludes many $a$.

\emph{Filter N (non-3-Higgs witness).}  If $m\mid a$ with $a/m$ odd, then
$2^m+1\mid 2^a+1$, so every prime factor of $2^m+1$ divides $n$.  Under H3, a
candidate is excluded as soon as some $2^m+1$ of this kind has a prime factor
that is not 3-Higgs.  Factorizations come from the bundled cache and
FactorDB~\cite{FactorDB}, and a partial factorization suffices when it exhibits
such a prime.  For example $a=4527$ in the class of $3^4\,41$ is excluded via
$m=1509$: the prime $20127043$ divides $2^{1509}+1$, and
$20127043-1=2\cdot3^4\cdot13\cdot19\cdot503$ has $v_3=4>3$.

\emph{Filter O (2-adic overshoot) is withdrawn.}  Version~2 closed the prime
powers of $K$ and of the seed under the rule ``$r\mid p^e+1$ forces $r$'', raising
exponents as incoming valuations grew, and stopped when the running sum of
$v_2(p^e+1)$ exceeded $a+1$.  It claimed this sum was a lower bound for the
left side of~\eqref{eq:budget}.  It is not: when the true exponent of $p$ is
larger than the one used, $v_2(p^e+1)$ can be smaller (it equals $1$ for every
even $e$), and $p^e+1$ need not divide the true $p^{e'}+1$, so neither the 2-adic
sum nor the forced primes are justified.

\subsection{Results}

\begin{proposition}\label{prop:ten}
No unitary perfect number refines the impostor kernel $K$ with $2^a\,\Vert\, n$
for any of the ten candidates that version~2 closed with filter O:
\[
(3^2\,5^3;\ a=10,30),\quad (3^4\,41;\ a=9,45,63,117),\quad
(5^2\,13^2;\ a=6,18,30),\quad (5^4\,29\,157^2\,313;\ a=26).
\]
This holds unconditionally.
\end{proposition}

\begin{proof}
Run the search of Section~\ref{sec:search} with the prime powers of $K$
placed as components before the search starts.  The other components are
taken in increasing order; in Lemma~\ref{lem:prefix}, $c_j$ then denotes the
largest of them chosen so far ($c_0=1$), and the forced components are counted
in $F_j$, $S_j$, $N_j$ and $D_j$.  For $(5^2\,13^2;\,a=30)$
nothing needs to be searched: $5^2\,\Vert\, 2^{30}+1$ and $13^2+1=2\cdot5\cdot17$,
so~\eqref{eq:exponent} would give $v_5(m)\ge3$, contradicting $5^2\,\Vert\, m$.
The other nine searches end with no solution and no open branch; they visit
$3$, $22$, $2$, $366$, $2737$, $690726$, $3$, $54$ and $212$ nodes respectively
(in the order listed, omitting $(5^2\,13^2;30)$).  The case $a=117$ was run
with the Python program, because $2^{117}+1$ has the $22$-digit prime factor
$5302306226370307681801$; this is the only prime above $2^{64}$ that the run uses,
and its primality was proved with PARI/GP.  As a positive control, the
same program with $5^4$ placed first at $a=18$ finds Wall's number.
\end{proof}

\begin{theorem}\label{thm:conditional}
Assume Hypothesis~H3.  No unitary perfect number refines a kernel containing
$13^2$ or $157^2$.  For $K\in\{3^2\,5^3,\ 3^4\,41\}$ and every $a$ in its seed
class with $1\le a\le 10000$, no unitary perfect number with $2^a\,\Vert\, n$
refines $K$.
\end{theorem}

\begin{proof}
The prime $17$ is not 3-Higgs, since $17-1=2^4$, and it divides both
$13^2+1=2\cdot5\cdot17$ and $157^2+1=2\cdot5^2\cdot17\cdot29$; by
\eqref{eq:exponent} it would divide $n$.  This disposes of three of the five
kernels for every $a$, which accounts for $1063$ of the $2119$ candidates with
$a\le10000$.  For the other two kernels there are $1056$ candidates: filter Z
excludes $262$ of them and filter N excludes $788$, both using H3, and the
remaining six are candidates of Proposition~\ref{prop:ten}.
\end{proof}

Theorem~\ref{thm:conditional} replaces version~2's Theorem~3 and Corollary~4.
It is conditional on H3, and filter O is replaced by
Proposition~\ref{prop:ten}.  The candidates with
$a$ covered by Theorem~\ref{thm:intro} are also excluded unconditionally.

\section{The set \texorpdfstring{$H_{\mathrm{even}}$}{H even}}\label{sec:Heven}

The results of this section are unconditional statements about the numbers
$2^m+1$ and the set $\mathcal P_3$ of 3-Higgs primes.  They bear on unitary
perfect numbers only through Hypothesis~H3: under H3 every prime factor of the
seed $2^a+1$ is 3-Higgs, so $a\in H$, and the same holds for every $m\mid a$
with $a/m$ odd.  Version~2 presented this connection as unconditional.

Let $H := \{ m \ge 1 : \text{every prime factor of } 2^m + 1
\text{ is 3-Higgs} \}$ and $H_{\mathrm{even}} := H \cap 2\mathbb{Z}$.

\begin{proposition}[Structural lemma]\label{prop:Heven-structure}
Suppose $m = 2k \in H_{\mathrm{even}}$ with $k$ odd and $k \ge 1$. Then
\begin{enumerate}[label=(\arabic*),leftmargin=*]
\item Every prime factor of $k$ is a 3-Higgs prime.
\item For every prime $q \mid k$, $v_q(k) \le 3$.
\item For every odd divisor $d \mid k$, $2d \in H_{\mathrm{even}}$.
\end{enumerate}
The first two conditions say $k$ is \emph{Higgs-bounded}: a product
of 3-Higgs primes each with multiplicity at most $3$. The third
condition is a multiplicative closure: writing
$S = \{k \text{ odd} : 2k \in H_{\mathrm{even}}\}$, $S$ is closed
under taking odd divisors.
\end{proposition}

\begin{proof}
By Zsigmondy's theorem applied to $\{2^n+1\}_{n \ge 1}$, the number
$2^{2k}+1$ has a primitive prime divisor $r$ for every $k \ge 1$ odd
(the only Zsigmondy exception for the sequence $2^n+1$ is $n = 3$,
which does not arise since $2k$ is even). The order of $2$ modulo $r$
is exactly $4k$, so $r \equiv 1 \pmod{4k}$ and $r - 1 = 4k \cdot s$
for some positive integer $s$.

By hypothesis $r$ is 3-Higgs, so every prime $q \mid r - 1$ is
3-Higgs and $v_q(r-1) \le 3$.

\noindent (1) For each prime $q \mid k$, $q \mid 4k \mid r - 1$, so $q$ is
3-Higgs.

\noindent (2) For each odd prime $q \mid k$,
$v_q(r-1) = v_q(4k) + v_q(s) = v_q(k) + v_q(s) \ge v_q(k)$. Combined
with $v_q(r-1) \le 3$, we get $v_q(k) \le 3$. (The prime $2$ does not
constrain $k$ here since $k$ is odd.)

\noindent (3) If $d \mid k$ with $d$ odd and $d < k$, then
$2d \mid 2k$ with $(2k)/(2d) = k/d$ odd. Therefore
$2^{2d}+1 \mid 2^{2k}+1$, so every prime divisor of $2^{2d}+1$ is a
prime divisor of $2^{2k}+1$ and hence 3-Higgs. So $2d \in H_{\mathrm{even}}$.
\end{proof}

\begin{proposition}[Fermat-prime obstruction]\label{prop:Heven-dyadic}
$H_{\mathrm{even}} \subseteq \{ m \equiv 2 \pmod 4 \}$.
\end{proposition}

\begin{proof}
Write $m = 2^j \cdot \ell$ with $\ell$ odd. We show that if $j \ge 2$
then $m \notin H$, which gives the conclusion.

Fix $k \ge 2$ and let $q$ be any prime divisor of the Fermat number
$F_k = 2^{2^k} + 1$. The order of $2$ modulo $q$ divides $2^{k+1}$ and
does not divide $2^k$ (since $q \mid F_k$ means $2^{2^k} \equiv -1
\pmod q$), so $\mathrm{ord}_q(2) = 2^{k+1}$. Consequently
\[
q \equiv 1 \pmod{2^{k+1}}.
\]
The refinement due to Lucas strengthens this for $k \ge 2$ to
$q \equiv 1 \pmod{2^{k+2}}$, hence
\[
v_2(q - 1) \ge k + 2 \ge 4 > 3 \quad (k \ge 2).
\]
The 3-Higgs definition forbids any prime factor of $p - 1$ to occur
with exponent exceeding $3$, so $q$ is not 3-Higgs.

Now $q \mid 2^m + 1$ iff $\mathrm{ord}_q(2) \mid 2m$ and
$\mathrm{ord}_q(2) \nmid m$, i.e., $v_2(\mathrm{ord}_q(2)) = v_2(m) + 1$.
With $\mathrm{ord}_q(2) = 2^{k+1}$ this is $v_2(m) = k$.

Therefore: for any $m$ with $v_2(m) = j \ge 2$, take $k = j$ and any
prime divisor $q$ of $F_j$; then $q \mid 2^m + 1$ and $q$ is not
3-Higgs, so $m \notin H$. The remaining cases $v_2(m) \le 1$ are
exactly $m$ odd ($v_2 = 0$) and $m \equiv 2 \pmod 4$ ($v_2 = 1$),
giving the stated containment for the even part of $H$.
\end{proof}

\begin{conjecture}\label{conj:Heven-finite}
$H_{\mathrm{even}}$ is finite.
\end{conjecture}

\begin{theorem}[Reduction to the prime case]\label{thm:prime-reduction}
$H_{\mathrm{even}}$ is finite if and only if
\[
H_{\mathrm{even}}^{\mathrm{prime}} \;:=\;
\{m = 2p \in H_{\mathrm{even}} : p \text{ odd prime}\}
\]
is finite. More precisely, if $H_{\mathrm{even}}^{\mathrm{prime}}$
has cardinality $N$ (with corresponding primes $p_1, \dots, p_N$),
then
\[
|H_{\mathrm{even}}| \;\le\; \prod_{i=1}^{N} 4 \;=\; 4^{N}.
\]
\end{theorem}

\begin{proof}
The forward direction is trivial. For the converse, let $m = 2k \in
H_{\mathrm{even}}$ with $k \ge 1$ odd. By
Proposition~\ref{prop:Heven-structure}(1)-(3), $k$ is Higgs-bounded
and every odd divisor $d \mid k$ satisfies $2d \in H_{\mathrm{even}}$.
In particular, for every prime $q \mid k$, $2q \in
H_{\mathrm{even}}^{\mathrm{prime}}$, so $q \in
\{p_1, \dots, p_N\}$. Combined with $v_q(k) \le 3$, this forces
$k$ to be of the form $\prod_{i=1}^N p_i^{e_i}$ with each $e_i \in
\{0,1,2,3\}$, giving at most $4^N$ possible $k$.
\end{proof}

\noindent The reduction is sharp in the sense that for every
$k = \prod p_i^{e_i}$ with each $2p_i \in H_{\mathrm{even}}^{\mathrm{prime}}$
and $e_i \le 3$, the candidate $m = 2k$ \emph{passes} the structural
lemma but may still fail (some prime factor of $2^{2k}+1$ may be
non-3-Higgs even though each prime in $k$ is 3-Higgs). So $4^N$ is a
candidate bound, refined to the actual $H_{\mathrm{even}}$ by checking
each candidate.

\smallskip
Theorem~\ref{thm:prime-reduction} reduces
Conjecture~\ref{conj:Heven-finite} to the prime case:
\emph{prove that only finitely many odd primes $p$ have all prime
factors of $2^{2p}+1$ in $\mathcal{P}_3$}. The remaining open candidates in Theorems~\ref{thm:Heven-3000}--\ref{thm:Heven-50000}
are controlled by the same prime-branch obstruction: the primitive
undecided cases have the form $m=2p$, while five composite-$k$
candidates ($m \in \{27978, 30354, 31538, 41898, 46630\}$, each with
$k$ a product of two Higgs primes) are inherited from unresolved
prime divisors through Proposition~\ref{prop:Heven-structure}.

\begin{theorem}\label{thm:Heven-1200}
$H_{\mathrm{even}} \cap [2, 1200] = \{ 2, 6, 10, 18, 26, 30, 46, 62, 82, 122 \}$.
\end{theorem}

\begin{proof}
By Proposition~\ref{prop:Heven-structure}, any $m = 2k \in
H_{\mathrm{even}}$ with $k$ odd and $k \le 600$ requires $k$ to be
Higgs-bounded. Of the $300$ odd $k$ in $[1, 600]$, exactly $246$ are
Higgs-bounded; the other $54$ are structurally excluded from
$H_{\mathrm{even}}$ without any factoring required.

For each of the $246$ Higgs-bounded candidates, we factor $2^m + 1$
(local Pollard~$\rho$/$p-1$ for small $m$, FactorDB~\cite{FactorDB}
for larger $m$) and recursively verify the 3-Higgs status of every
listed prime factor.
Of the $246$ candidates:
\begin{itemize}
\item $10$ values of $m$ have every prime factor of $2^m+1$ 3-Higgs;
these are exactly the elements listed in the theorem statement.
\item $236$ values of $m$ each contain at least one verified
non-3-Higgs prime factor (witness digit sizes ranging from $3$ to
$115$).
\item $0$ values remain undecided.
\end{itemize}
The script \texttt{verify\_Heven\_rigorous.py} is provided in the
supplementary materials and records, for every $m$ in the second
group, the specific non-3-Higgs prime witness.
\end{proof}

Theorem~\ref{thm:Heven-1200} extends the verified gap considerably:
$H_{\mathrm{even}}$ has no element in the wide range $(122, 1200]$,
strongly supporting Conjecture~\ref{conj:Heven-finite}.

\begin{theorem}\label{thm:Heven-3000}
$H_{\mathrm{even}} \cap (1200, 3000]$ is contained in the two-element
candidate set $\{2426, 2602\}$. Every other even $m$ in $(1200, 3000]$
is rigorously excluded from $H_{\mathrm{even}}$, by either the
structural lemma (Proposition~\ref{prop:Heven-structure}), an
explicit non-3-Higgs prime witness in the factorization of $2^m + 1$,
or the deep-Pratt descent of Lemma~\ref{lem:closures-deep} (which
closes $m = 2446$).
\end{theorem}

\begin{proof}
Of the $450$ odd $k$ in $[601, 1500]$, $111$ are not Higgs-bounded
and excluded by Proposition~\ref{prop:Heven-structure}. The remaining
$339$ candidates split: $336$ have a verified non-3-Higgs prime
witness directly in $2^m + 1$ (recorded by
\texttt{verify\_Heven\_rigorous.py}); $3$ were initially blocked by
unfactored composite cofactors. One of these three, $m = 2446$, is
closed by Lemma~\ref{lem:closures-deep} (a $368$-digit
APR-CL-verified prime divisor of $2^{2446}+1$ has $4513 \mid p^{*}-1$,
and $v_2(4512) = 5 > 3$). The remaining two $m$ values, $\{2426, 2602\}$,
have FactorDB factorizations whose composite cofactors are not yet
deep enough for a Pratt-tree closure.

For $m = 1202$, the cache records an $87$-digit prime factor
$P_1 \mid 2^{1202}+1$ with
$P_1 - 1 = 4 \cdot 601 \cdot Q$, where $Q$ is a prime. Moreover
$Q - 1 = 2^{10} \cdot 3 \cdot 337 \cdot R$ for a $78$-digit cofactor
$R$ (the explicit small-prime decomposition is verified by trial
division). Since $v_2(Q-1) = 10 > 3$, the prime $Q$ fails the 3-Higgs
$v_2$ exponent bound. Consequently $P_1$ has a non-3-Higgs prime
factor $Q$ in $P_1 - 1$ (with $v_Q(P_1-1) = 1$), so $P_1$ is itself
not 3-Higgs. Since $P_1 \mid 2^{1202}+1$, the seed $2^{1202}+1$
contains a non-3-Higgs prime, and $m = 1202 \notin H_{\mathrm{even}}$.
(The cache also records a $91$-digit prime factor $P$ of $2^{1202}+1$ with
$v_2(P-1)=6$, which is a witness by itself.)

For $m = 2426$ and $m = 2602$, the FactorDB factorizations of
$2^m + 1$ contain composite cofactors of $355$--$392$ digits
(unfactored), so we cannot currently rule out hidden non-3-Higgs
factors.

All other Higgs-bounded candidates have at least one explicit
non-3-Higgs prime witness recorded by
\texttt{verify\_Heven\_rigorous.py}.
\end{proof}

\noindent The two open candidates $\{2426, 2602\}$ would
resolve with NFS-level factoring of $2^m + 1$ for those $m$. The
result $|H_{\mathrm{even}} \cap [2, 3000]| \le 12$ is striking
evidence for Conjecture~\ref{conj:Heven-finite}.

\begin{theorem}\label{thm:Heven-5000}
$H_{\mathrm{even}} \cap (3000, 5000]$ is contained in the five-element
candidate set
\[
\{3398, 3518, 4166, 4502, 4622\}.
\]
\end{theorem}

\begin{proof}
Of $500$ odd $k$ in $[1501, 2500]$, $132$ are not Higgs-bounded and
excluded by Proposition~\ref{prop:Heven-structure}. The remaining
$368$ candidates split: $363$ values of $m = 2k$ have a verified
non-3-Higgs prime witness in $2^m + 1$, and $5$ values are blocked by
unfactored composite cofactors in FactorDB's partial factorization
(listed in the theorem). Examples of closures: $m = 4366$ is killed
by $593 \mid 2^{4366}+1$ ($593 - 1 = 2^4 \cdot 37$, $v_2 = 4$);
$m = 4742$ by $493169 \mid 2^{4742}+1$
($493169 - 1 = 2^4 \cdot 30823$); $m = 4022$ by a recursive Pratt-tree
chain: the 18-digit prime
$920793289987614829 \mid 2^{4022}+1$ has its $(p-1)$ chain descend
through $1230855683509 \to 2169031 \to 17$, where the final factor
$17$ is non-3-Higgs (since $17-1 = 2^4$). All $363$ rigorous exclusions
are recorded by the supplementary script.
\end{proof}

\begin{theorem}\label{thm:Heven-10000}
$H_{\mathrm{even}} \cap (5000, 10000]$ is contained in the
five-element candidate set
\[
\{5114, 5774, 7846, 8966, 9326\}.
\]
\end{theorem}

\begin{proof}
Of $1250$ odd $k$ in $[2501, 5000]$, $364$ are not Higgs-bounded
and excluded by Proposition~\ref{prop:Heven-structure}. The remaining
$886$ Higgs-bounded candidates split as $881$ verified non-3-Higgs
exclusions (via small-prime witnesses in $2^m + 1$ recorded by the
supplementary script) and $5$ partial-cofactor unknowns (the five $m$
values listed in the theorem, each blocked by an unfactored composite
cofactor of 1500--2800 digits in $2^m+1$).
\end{proof}

\begin{theorem}\label{thm:Heven-15000}
$H_{\mathrm{even}} \cap (10000, 15000]$ is contained in the
ten-element candidate set
\[
\{10118, 10454, 11062, 11794, 11878, 12778, 13382, 13966, 14642, 14698\}.
\]
\end{theorem}

\begin{proof}
Of $1250$ odd $k$ in $[5001, 7500]$, $383$ are not Higgs-bounded
and excluded by Proposition~\ref{prop:Heven-structure}. The remaining
$867$ Higgs-bounded candidates yield $m = 2k$ values for which
FactorDB partial factorizations provide small-prime data. Of these,
exactly $857$
have a verified non-3-Higgs prime witness (recorded by the
supplementary script \texttt{verify\_Heven\_rigorous.py}, or, for
$m=10294$ and $m=10958$, given by Lemma~\ref{lem:closures-deep}), and $10$
are partial-cofactor unknowns listed in the theorem.  In this range
these are primitive candidates $m=2p$, with
$p$ an odd prime in $\{5059, 5227, 5531, 5897, 5939, 6389, 6691,
6983, 7321, 7349\}$, and each is blocked by one or more
unfactored composite cofactors in $2^m+1$ of $1300$--$4500$ digits.
For all ten, we have additionally queried FactorDB for the
Aurifeuillean halves $L_p$ and $M_p$ (Equation~\ref{eq:Aurifeuillean});
no new small prime factor in any $L_p$ or $M_p$ failed the 3-Higgs
test, so the obstruction is genuinely NFS-blocked rather than
shallow. (Two values $m = 10294$ and $m = 10958$, originally in this
range as partial-cofactor unknowns, were closed via deep Pratt-tree
descent on $1549$- and $1649$-digit APR-CL-verified prime divisors of
$2^m + 1$; see Lemma~\ref{lem:closures-deep}.)
\end{proof}

\begin{theorem}\label{thm:Heven-20000}
$H_{\mathrm{even}} \cap (15000, 20000]$ is contained in the
twenty-seven-element candidate set
\[
\begin{aligned}
\{ & 15122, 15182, 15206, 15214, 15754, 15758, 16358, 16574, 16778,\\
   & 16838, 16922, 17086, 17126, 17162, 17338, 17726, 18134,\\
   & 18418, 18934, 18958, 19078, 19226, 19322, 19718, 19802,\\
   & 19846, 19862 \}.
\end{aligned}
\]
\end{theorem}

\begin{proof}
Of $1250$ odd $k$ in $[7501, 10000]$, $392$ are not Higgs-bounded
and excluded by Proposition~\ref{prop:Heven-structure}. The remaining
$858$ Higgs-bounded candidates yield $m = 2k$ values for which
FactorDB partial factorizations provide small-prime data. Of these,
exactly $831$ have a verified non-3-Higgs prime witness; the
remaining $27$ are partial-cofactor unknowns listed in the theorem.
In this range these are primitive candidates $m=2p$ with $p$ an odd
Higgs prime in $[7561, 9931]$, blocked by unfactored composite
cofactors of $1700$--$6000$ digits. (Two values
$m = 17398$ and $m = 19066$, which appeared as partial-cofactor
unknowns in our earlier verification, were subsequently closed via
non-3-Higgs Pratt-tree witnesses extracted from $p-1$ for large
APR-CL-verified prime divisors $p \mid 2^m + 1$; see Lemma~\ref{lem:closures-deep}.)
For all twenty-seven remaining, the Aurifeuillean halves $L_p, M_p$
(Equation~\ref{eq:Aurifeuillean}) were queried in FactorDB; no
new small prime factor failed the 3-Higgs test.
\end{proof}

\begin{theorem}\label{thm:Heven-25000}
$H_{\mathrm{even}} \cap (20000, 25000]$ is contained in the
twenty-seven-element candidate set
\[
\begin{aligned}
\{ & 20138, 20338, 20506, 20662, 20926, 20974, 21118, 21302,\\
   & 21334, 21466, 21818, 21958, 22054, 22234, 22262, 22394, 22486,\\
   & 22642, 22706, 23194, 23402, 23578, 23878, 23942, 24082, 24298,\\
   & 24914 \}.
\end{aligned}
\]
\end{theorem}

\begin{proof}
Of $1250$ odd $k$ in $[10001, 12500]$, $416$ are not Higgs-bounded
and excluded by Proposition~\ref{prop:Heven-structure}. The remaining
$834$ Higgs-bounded candidates yield $m = 2k$ values with FactorDB
partial factorizations; of these, $807$ have a verified non-3-Higgs
prime witness in $2^m + 1$, and $27$ are partial-cofactor unknowns
listed in the theorem.  In this range these are primitive candidates
$m=2p$ with $p$ an odd prime in $[10069, 12457]$.
(The value $m = 20282$ was closed via a Pratt-tree witness $v_2(p-1) = 5071$
for a $3053$-digit APR-CL-verified prime divisor of $2^{20282}+1$; see
Lemma~\ref{lem:closures-deep}.)
\end{proof}

\begin{theorem}\label{thm:Heven-30000}
$H_{\mathrm{even}} \cap (25000, 30000]$ is contained in the
forty-eight-element candidate set
\[
\begin{aligned}
\{ & 25022, 25082, 25106, 25294, 25486, 25526, 25646, 25778, 25786,\\
   & 25822, 26066, 26098, 26198, 26302, 26342, 26438, 26482, 26534,\\
   & 26618, 26654, 26662, 26678, 26798, 26902, 27374, 27446, 27526,\\
   & 27662, 27806, 27926, 27978, 27998, 28058, 28102, 28174, 28214,\\
   & 28318, 28442, 28586, 28654, 28862, 29098, 29114, 29126, 29558,\\
   & 29642, 29662, 29914 \}.
\end{aligned}
\]
\end{theorem}

\begin{proof}
Of $1250$ odd $k$ in $[12501, 15000]$, $427$ are not Higgs-bounded
and excluded by Proposition~\ref{prop:Heven-structure}. The remaining
$823$ Higgs-bounded candidates yield $m = 2k$ values; of these,
$775$ have a verified non-3-Higgs prime witness in $2^m + 1$,
and $48$ are partial-cofactor unknowns listed in the theorem; the primitive cases in this list are of the form $m=2p$.
\end{proof}

\begin{theorem}\label{thm:Heven-35000}
$H_{\mathrm{even}} \cap (30000, 35000]$ is contained in the
thirty-two-element candidate set
\[
\begin{aligned}
\{ & 30354, 30386, 30542, 30878, 30994, 31162, 31258, 31454,\\
   & 31466, 31538, 31634, 31802, 31918, 31942, 31982, 32126, 32174,\\
   & 32534, 32762, 32902, 33134, 33214, 33382, 33386, 33574, 33778,\\
   & 33806, 33974, 34058, 34246, 34834, 34934 \}.
\end{aligned}
\]
\end{theorem}

\begin{proof}
Of $1250$ odd $k$ in $[15001, 17500]$, $429$ are not Higgs-bounded
and excluded by Proposition~\ref{prop:Heven-structure}. The remaining
$821$ Higgs-bounded candidates yield $m = 2k$ values; of these,
$789$ have a verified non-3-Higgs prime witness in $2^m + 1$, and
$32$ are partial-cofactor unknowns listed in the theorem; the primitive cases in this list are of the form $m=2p$. (The value $m = 30882$ was closed
via Proposition~\ref{prop:Heven-structure}(3) from the closure of
$m = 10294$; see Lemma~\ref{lem:closures-deep}.)
\end{proof}

\begin{theorem}\label{thm:Heven-40000}
$H_{\mathrm{even}} \cap (35000, 40000]$ is contained in the
thirty-five-element candidate set
\[
\begin{aligned}
\{ & 35038, 35194, 35366, 35654, 35678, 35818, 35858, 35878, 35942,\\
   & 35978, 36094, 36122, 36298, 36458, 36574, 36602, 36682, 36794,\\
   & 36986, 37234, 37322, 37358, 37382, 38174, 38282, 38518, 38954,\\
   & 39014, 39062, 39154, 39502, 39706, 39826, 39926, 39958 \}.
\end{aligned}
\]
\end{theorem}

\begin{proof}
Of $1250$ odd $k$ in $[17501, 20000]$, $419$ are not Higgs-bounded
and excluded by Proposition~\ref{prop:Heven-structure}. The remaining
$831$ Higgs-bounded candidates yield $m = 2k$ values; of these,
$796$ have a verified non-3-Higgs prime witness in $2^m + 1$, and
$35$ are partial-cofactor unknowns listed in the theorem; the primitive cases in this list are of the form $m=2p$.
\end{proof}

\begin{theorem}\label{thm:Heven-45000}
$H_{\mathrm{even}} \cap (40000, 45000]$ is contained in the
thirty-three-element partial-cofactor candidate set
\[
\begin{aligned}
\{ & 40234, 40466, 40538, 40574, 40822, 40886, 41126, 41326, 41386,\\
   & 41714, 41758, 41878, 41898, 42134, 42314, 42358, 42646, 42766,\\
   & 43058, 43154, 43322, 43402, 43574, 43634, 43642, 43678, 43786,\\
   & 44054, 44186, 44294, 44378, 44518, 44582 \}.
\end{aligned}
\]
\end{theorem}

\begin{proof}
Of the $1250$ odd $k$ in $[20001, 22500]$, $428$ are not
Higgs-bounded and excluded by Proposition~\ref{prop:Heven-structure}.
The remaining $822$ Higgs-bounded candidates yield $m = 2k$ values;
of these, $789$ have a verified non-3-Higgs prime witness in
$2^m + 1$ (after a second-pass FactorDB retry that recovered
factor data for $26$ initially-no-factor cases), and $33$ are
partial-cofactor unknowns listed in the theorem; the primitive cases in this list are of the form $m=2p$.
\end{proof}

\begin{theorem}\label{thm:Heven-50000}
$H_{\mathrm{even}} \cap (45000, 50000]$ is contained in the
thirty-eight-element partial-cofactor candidate set
\[
\begin{aligned}
\{ & 45142, 45286, 45394, 45434, 45478, 45566, 45622, 45706, 45842,\\
   & 46022, 46042, 46058, 46394, 46454, 46630, 46714, 46862, 47206,\\
   & 47246, 47342, 47378, 47438, 47506, 47858, 47962, 47986, 48086,\\
   & 48218, 48338, 48362, 48838, 49354, 49466, 49694, 49834, 49906,\\
   & 49958, 49978 \}.
\end{aligned}
\]
\end{theorem}

\begin{proof}
Of the $1250$ odd $k$ in $[22501, 25000]$, $430$ are not
Higgs-bounded and excluded by Proposition~\ref{prop:Heven-structure}.
The remaining $820$ Higgs-bounded candidates yield $m = 2k$ values;
of these, $782$ have a verified non-3-Higgs prime witness in
$2^m + 1$ (after a second-pass FactorDB retry that recovered
factor data for $18$ initially-no-factor cases), and $38$ are
partial-cofactor unknowns listed in the theorem; the primitive cases in this list are of the form $m=2p$.
\end{proof}

Without the seven closures of Lemma~\ref{lem:closures-deep}, the
factor-cache verification alone gives
\[
|H_{\mathrm{even}} \cap [2, 50000]| \le 279,
\]
with the 10 verified elements through $m = 122$, no verified element
in $(122, 50000]$, and at most $269$ undecided candidates beyond
($198$ in $(1200, 40000]$, $33$ in $(40000, 45000]$ and $38$ in
$(45000, 50000]$). The seven closures, all of whose witness primes are
APR-CL verified (see the proof of Lemma~\ref{lem:closures-deep}), remove
seven candidates in $(1200, 40000]$. They are already applied in
Theorems~\ref{thm:Heven-3000}--\ref{thm:Heven-50000}, which together with
Theorem~\ref{thm:Heven-1200} give
\[
|H_{\mathrm{even}} \cap [2, 50000]| \le 272
\qquad
\text{(see the note on primality below),}
\]
with $262$ undecided candidates remaining (the $71$ in $(40000, 50000]$
admit no Pratt-tree closures in the current cache; their resolution
requires further FactorDB data or local ECM at higher $B_1$).  The scripts test primality with a 12-base Miller--Rabin test; the large
primes of Lemma~\ref{lem:closures-deep} are proved prime by APR-CL.

\begin{lemma}[Deep-Pratt closures via large APR-CL-verified prime divisors]
\label{lem:closures-deep}
For each of the following six $m = 2p$ values, the indicated prime
$p^{*} \mid 2^m + 1$ has a non-3-Higgs witness in $p^{*} - 1$,
hence $p^{*} \notin \mathcal{P}_3$ and $m \notin H_{\mathrm{even}}$:
\begin{center}
\begin{tabular}{rrll}
\toprule
$m$ & $p^{*}$ digits & Non-3-Higgs witness $q$ & Reason \\
\midrule
$2446$  & $368$  & $q = 4513 \mid p^{*}-1$ & $v_2(4512) = 5 > 3$ \\
$10294$ & $1549$ & $q = 2657 \mid p^{*}-1$ & $v_2(2656) = 5 > 3$ \\
$10958$ & $1649$ & $q = 593 \mid p^{*}-1$ & $v_2(592) = 4 > 3$ \\
$17398$ & $2612$ & $q = 139313 \mid p^{*}-1$ & $v_2(139312) = 4 > 3$ \\
$19066$ & $2870$ & $q = 343081 \mid p^{*}-1$ & contains $17$ via $953$ \\
$20282$ & $3053$ & $v_2(p^{*}-1) = 5071$ & direct $v_2$ overflow \\
\bottomrule
\end{tabular}
\end{center}
In addition, $m = 30882 = 6 \cdot 5147$ is closed by
Proposition~\ref{prop:Heven-structure}(3) from the closure of
$m = 10294 = 2 \cdot 5147$.
\end{lemma}

\begin{proof}
For each row, the prime $p^{*}$ is recorded in the
supplementary factor cache and verified to divide $2^m + 1$ by direct
modular reduction. The primality of each $p^{*}$ is established by an
APR-CL primality verification carried out in PARI/GP 2.17.3
(\texttt{isprime(n, 2)}), and the subsequent exclusions use Pratt-tree
descents~\cite{FKL,WimmerNoschinski}. PARI's \texttt{isprime(n, 2)}
returns a Boolean rather than an independently checkable certificate
object;
the supplementary material therefore bundles the deterministic input
files (the $p^{*}$ values in decimal) and the runtime logs ($3.4$,
$79$, $88$, $615$, $679$, $958$ seconds wall-clock for rows
$1$--$6$ on $11$-thread parallel execution), so readers can re-run
\texttt{isprime(n, 2)} on the recorded version and reproduce the
result. Conditional on the correctness of PARI/GP's APR-CL
implementation, the primality of each $p^{*}$ is thereby established. For each row's witness $q$, trial division
verifies $q \mid p^{*} - 1$, and the non-3-Higgs status of $q$ follows
either from the explicit $v_2$ overflow (rows 1, 2, 3, 4, 6) or from a
Pratt-tree descent terminating at $17$ (row 5: $343081 \succ 953
\succ 17$, with $v_2(16) = 4 > 3$). Hence $p^{*}$ has a non-3-Higgs
prime in its Pratt tree and $p^{*} \notin \mathcal{P}_3$; since
$p^{*} \mid 2^m+1$, this gives $m \notin H_{\mathrm{even}}$. The
final case $m = 30882$ follows from
Proposition~\ref{prop:Heven-structure}(3): $10294 \mid 30882$, so
$2^{10294}+1 \mid 2^{30882}+1$; any non-3-Higgs prime of $2^{10294}+1$
(in particular the $1549$-digit $p^{*}$) is also a divisor of
$2^{30882}+1$.
\end{proof}

\subsection{Ford's theorem and thinness}

The structural lemma (Proposition~\ref{prop:Heven-structure}) confines
$H_{\mathrm{even}}$ to the doubled image of a much more restrictive
arithmetic object: the multiplicative semigroup
\[
\mathcal{S}_3^{(\le 3)}
\;:=\;
\left\{\, n \ge 1 \;:\;
\begin{array}{l}
\text{every prime factor of } n \text{ lies in } \mathcal{P}_3,\\
\text{and each such prime has exponent at most } 3
\end{array}
\right\}.
\]
where $\mathcal{P}_3$ denotes the 3-Higgs primes (OEIS A057447). The
defining recursive closure of $\mathcal{P}_3$,
\[
p \in \mathcal{P}_3 \implies \bigl(q \mid (p-1) \Rightarrow q \in \mathcal{P}_3\bigr),
\]
makes $\mathcal{P}_3$ a downward-closed set of primes in the sense of
Ford~\cite{Ford2014}. Since the smallest omitted prime is $p_0 = 17$
(as $17 - 1 = 2^4$ violates the $v_2 \le 3$ exponent bound), Ford's
theorem applies unconditionally and yields
\begin{equation}\label{eq:ford-bound}
\Pi_3(X) := \#\{p \le X : p \in \mathcal{P}_3\}
\;\ll\; \frac{X}{(\log X)^{1 + 1/(p_0 - 1)}}
\;=\; \frac{X}{(\log X)^{17/16}},
\end{equation}
together with $\Pi_3(X) \ll X^{1-\delta}$ for some absolute
$\delta > 0$, and the absolute convergence
$\sum_{p \in \mathcal{P}_3} 1/p < \infty$.

An entirely analogous argument applies to $H_{\mathrm{odd}} := H
\setminus H_{\mathrm{even}}$ (the odd part of $H$). For $m$ odd in
$H$ with $m > 3$, the Zsigmondy theorem yields a primitive prime
divisor $r$ of $2^{2m}-1$ with $r \mid 2^m+1$; this $r$ has
$\operatorname{ord}_r(2) = 2m$ and so $2m \mid r-1$, giving the
same recursive constraint on the prime factors of $m$ themselves.
The exceptional small values $m = 1$ ($2^1+1 = 3 \in \mathcal{P}_3$,
$1$ is the empty product) and $m = 3$ ($2^3+1 = 9 = 3^2$, where the
Zsigmondy exception $2^6 - 1$ removes the primitive-divisor step but
$3 \in \mathcal{P}_3$ directly) are verified by hand. Thus
$m \in \mathcal{S}_3^{(\le 3)}$ as well.

\begin{theorem}\label{thm:H-thin}
There exists an absolute constant $\eta > 0$ such that
\[
\#\{m \le X : m \in H\} \;\ll\; X^{1-\eta}.
\]
Moreover, $\displaystyle\sum_{m \in H} \frac{1}{m} < \infty$.
The same bound holds for the subsets $H_{\mathrm{even}}$ and
$H_{\mathrm{odd}}$.
\end{theorem}

\begin{proof}
By the structural lemma (Proposition~\ref{prop:Heven-structure}) and
its analog above, every $m \in H$ satisfies $m \in \mathcal{S}_3^{(\le 3)}$
or $m/2 \in \mathcal{S}_3^{(\le 3)}$. So
$\#\{m \le X : m \in H\} \le 2 \cdot \#\{k \le X : k \in
\mathcal{S}_3^{(\le 3)}\}$.

For $s \in (1-\delta, 1)$,
\[
\#\{n \le X : n \in \mathcal{S}_3^{(\le 3)}\}
\;\le\; X^s \prod_{p \in \mathcal{P}_3} (1 + p^{-s} + p^{-2s} + p^{-3s})
\]
by Rankin's trick. The Euler product converges because $\sum_{p \in
\mathcal{P}_3} p^{-s}$ converges for any $s > 1-\delta$ (by partial
summation against (\ref{eq:ford-bound})). Hence
$\#\{n \le X : n \in \mathcal{S}_3^{(\le 3)}\} \ll X^s$, and choosing
$s = 1 - \eta$ for any $\eta < \delta$ gives the claim.

The convergence of $\sum 1/m$ follows from
$\sum_{n \in \mathcal{S}_3^{(\le 3)}} 1/n =
\prod_{p \in \mathcal{P}_3}(1 + p^{-1} + p^{-2} + p^{-3}) < \infty$.
\end{proof}

\begin{corollary}[Power-saving counting and reciprocal mass for the exponent-capped family]\label{cor:A-bound}
With $A(x) := \#(\mathcal{S}_3^{(\le 3)} \cap [1, x])$, there is an absolute
constant $\eta > 0$ (the same $\eta$ as in
Theorem~\ref{thm:H-thin}) such that
\[
A(x) \;\ll\; x^{1-\eta}
\qquad \text{and} \qquad
\sum_{n \in \mathcal{S}_3^{(\le 3)}} \frac{1}{n} < \infty.
\]
\end{corollary}

\begin{proof}
Both statements are the intermediate steps in the proof of
Theorem~\ref{thm:H-thin}: the Rankin-trick application together with
Ford's bound (\ref{eq:ford-bound}) and the Euler-product convergence.
\end{proof}

The corollary is qualitatively stronger than density zero: the set
$H_{\mathrm{even}}$ has a power-saving thin support, and its
reciprocals sum to a finite constant. The remaining gap between this
unconditional result and Conjecture~\ref{conj:Heven-finite} is the
last analytic step: showing that for every sufficiently large
$k \in \mathcal{S}_3^{(\le 3)}$, some prime $r \mid 2^{2k}+1$ with
$\operatorname{ord}_r(2) = 4k$ has $(r-1)/(4k) \notin \mathcal{S}_3^{(\le 3)}$.

\subsection{Aurifeuillean structure of the open candidates}\label{sec:Aurifeuillean}

Every primitive undecided candidate in
Theorems~\ref{thm:Heven-3000}--\ref{thm:Heven-50000}
has the form $m = 2p$ with $p$ an odd prime; composite undecided
candidates are inherited from unresolved prime divisors. For $m=2p$,
the classical Aurifeuillean
identity~\cite{Wagstaff2002} gives an algebraic factorization
\begin{equation}\label{eq:Aurifeuillean}
2^{2p} + 1
\;=\;
\bigl(2^p - 2^{(p+1)/2} + 1\bigr)
\bigl(2^p + 2^{(p+1)/2} + 1\bigr)
\;=:\; L_p \cdot M_p.
\end{equation}
Each $L_p$ and $M_p$ has approximately half the bit-length of
$2^{2p}+1$. Moreover, both sides are sparse polynomials of degree
four in a power of $2$:
\begin{itemize}
\item For $p = 4u+1$, set $X = 2^u$, so $2^p = 2X^4$ and
$2^{(p+1)/2} = 2X^2$. Then
$L_p = 2X^4 - 2X^2 + 1, \quad M_p = 2X^4 + 2X^2 + 1$.
\item For $p = 4u+3$, set $X = 2^u$, so $2^p = 8X^4$ and
$2^{(p+1)/2} = 4X^2$. Then
$L_p = 8X^4 - 4X^2 + 1, \quad M_p = 8X^4 + 4X^2 + 1$.
\end{itemize}
In both cases each half is an integer-coefficient quartic in $X$.
This is the natural starting point for SNFS polynomial selection;
In special-number field-sieve practice, polynomials of
degree $5$ or $6$ (obtained by re-parameterization) may sieve
more efficiently than the natural quartic for numbers in the
$300$--$3000$-digit range relevant here. The Aurifeuillean halves are special-form targets rather than
generic hard composites: each remaining prime-branch candidate gives
an SNFS target with explicit small-degree algebraic structure.

\paragraph{Worked example: $m = 2426$, $p = 1213$.}
For $p = 1213 \equiv 1 \pmod 4$, write $1213 = 4 \cdot 303 + 1$ and
set $x = 2^{303}$. Then
\[
2^{2426}+1 = (2x^4 - 2x^2 + 1)(2x^4 + 2x^2 + 1).
\]
Numerically, $L_{1213}$ and $M_{1213}$ are each $\approx 366$ decimal
digits. The known prime factor $P = 25893760589$ divides $L_{1213}$
exactly (not $M_{1213}$). A complete Pratt-tree descent~\cite{Ford2014}
on $P$, namely
\begin{align*}
P-1 &= 2^2 \cdot 1213 \cdot 5336719 \\
1213-1 &= 2^2 \cdot 3 \cdot 101, \qquad 101-1 = 2^2 \cdot 5^2 \\
5336719-1 &= 2 \cdot 3 \cdot 889453, \qquad 889453-1 = 2^2 \cdot 3^2 \cdot 31 \cdot 797 \\
797-1 &= 2^2 \cdot 199, \qquad 199-1 = 2 \cdot 3^2 \cdot 11, \qquad 11-1 = 2 \cdot 5,
\end{align*}
visits the distinct primes
$\{2,3,5,11,31,101,199,797,1213,889453,5336719,P\}$ with maximal
tree height $8$, and every exponent in every intermediate $q-1$ is at
most $3$. So $P$ is \emph{fully} 3-Higgs-compatible; no shallow
recursive obstruction exists. Closing $m = 2426$ therefore requires
finding a non-3-Higgs prime in the residual composite
$L_{1213} / (5 \cdot P)$ ($355$ digits) or in $M_{1213}$ ($366$
digits) --- both special-form SNFS targets after the split.

This example is independently confirmed by the official Cunningham
Project tables~\cite{Wagstaff2002}: the January 2026 base-$2$ table
lists
\[
2,\,2426\mathrm{L} \;=\; 5 \cdot 25893760589 \cdot \mathrm{C}355,
\qquad
2,\,2426\mathrm{M} \;=\; \mathrm{C}366,
\]
matching our decomposition exactly. The prime $25893760589$ is
therefore a tabulated Cunningham factor of an as-yet-unfinished
Aurifeuillean half, not an ad-hoc byproduct of our scripts.

\medskip
Under Hypothesis~H3, a proof of Conjecture~\ref{conj:Heven-finite}
would leave only finitely many
candidates in each even impostor class, for every $a$, after filter N; those
would still have to be settled by other means, such as the search of
Section~\ref{sec:search}.

\subsection{The reduction to \texorpdfstring{$H_{\mathrm{even}}$}{H even}}

Under Hypothesis~H3, the connection between $H_{\mathrm{even}}$ and the
impostor branch of Section~\ref{sec:higgs} is sharp:

\begin{proposition}\label{prop:Heven-reduction}
Let $K$ be an impostor kernel of even seed parity (i.e., the impostor
seed congruence has $a$ even). The set of $a$ in $K$'s seed congruence
class that survive filter $N$ (no proper divisor or self has a
non-3-Higgs prime in $2^m + 1$) is exactly
\[
\{a \in K\text{'s seed class} \cap H_{\mathrm{even}} :
   \text{every proper } m \mid a \text{ with } a/m \text{ odd is in } H_{\mathrm{even}}\}.
\]
\end{proposition}

\begin{proof}
Filter $N$ tests each $m$ in $\{m : m \mid a, a/m \text{ odd}\}$,
including $m = a$ itself, for the existence of a non-3-Higgs prime in
$2^m + 1$. By definition of $H_{\mathrm{even}}$, no such non-3-Higgs
prime exists iff every $m$ in this set lies in $H_{\mathrm{even}}$.
\end{proof}

\begin{corollary}[]\label{cor:finite-cascade}
If Conjecture~\ref{conj:Heven-finite} holds, then for each of the four
impostor kernels with even seed parity ($3^2\,5^3$, $5^2\,13^2$,
$5^4\,157^2\,313$, and $5^4\,29\,157^2\,313$), only finitely many
$a$ in the respective seed class survive filter $N$.
\end{corollary}

The fifth kernel $3^4\,41$ has odd seed parity ($a \equiv 9 \pmod{18}$);
its filter $N$ runs over odd $m$ and reduces to the analogous odd-parity
set $H_{\mathrm{odd}}$, for which an analogous finiteness conjecture
applies.

In particular, for the $3^2\,5^3$ kernel ($a \equiv 10 \pmod{20}$), the
intersection $H_{\mathrm{even}} \cap \{a : a \equiv 10 \pmod{20}\}$
consists, within the verified range, of exactly two elements:
$\{10, 30\}$. Both are eliminated unconditionally by
Proposition~\ref{prop:ten}. The other $248$ candidates in
$\{a \equiv 10 \pmod{20}, a \le 5000\}$ are excluded by filter $N$
or $Z$, under Hypothesis~H3.

\section{Open problems}\label{sec:open}

\begin{enumerate}[label=(\arabic*),leftmargin=*]
\item \textbf{Larger $a$.}  The search is complete for each $a$ it closes, and
its cost grows quickly with $a$; the obstruction is a prefix whose product is
already very close to $2^{a+1}/(2^a+1)$ while two or more components remain
unconstrained, so that the next component must be sought among all primes in a
long interval.  \RESULTopen{}  A method for these tails that does not scan the
interval, for example by fixing the exponents of the demanded primes first and
solving for the last two components by divisor enumeration, would extend
Theorem~\ref{thm:intro}.

\item \textbf{Hypothesis H3.}  Prove or refute that every prime divisor of a
unitary perfect number is 3-Higgs.  A proof would make
Theorem~\ref{thm:conditional} unconditional and would connect the
$H_{\mathrm{even}}$ problem of Section~\ref{sec:Heven} to unitary perfect
numbers.

\item \textbf{Finiteness of $H_{\mathrm{even}}$} (Conjecture~\ref{conj:Heven-finite}),
or more modestly a proof that only finitely many primes $p$ have every prime
factor of $2^{2p}+1$ in $\mathcal P_3$ (Theorem~\ref{thm:prime-reduction}).

\item \textbf{Formal verification.}  The search uses only exact arithmetic and
the lemmas of Section~\ref{sec:search}, so its correctness argument is short;
a certified version (for example, a proof checker for its node log) would remove
the need to trust the implementation.
\end{enumerate}

\section{Reproducibility}\label{sec:repro}

The ancillary directory \texttt{anc/} contains:
\begin{itemize}[leftmargin=*]
\item \texttt{search/search2.py}: the search of Section~\ref{sec:search} in
Python, with exact rational arithmetic (requires \texttt{sympy} and
\texttt{numpy});
\item \texttt{search/upn.c}: the same search in C with GMP, and
\texttt{search/run\_all.py}, which factors $2^a+1$, runs the search for each
$a$ of Theorem~\ref{thm:intro} and checks that no branch is open;
\item \texttt{search/logs/}: the output of those runs, and of the kernel runs of
Proposition~\ref{prop:ten};
\item the material of version~2 for Sections~\ref{sec:higgs}
and~\ref{sec:Heven}: the factor cache, the certificate and $H_{\mathrm{even}}$
scripts, and the APR-CL transcripts.
\end{itemize}
For example, the case $a=12$ is checked by
\begin{verbatim}
cc -O2 -o upn anc/search/upn.c -lgmp -lm
./upn 12 "17:1,241:1"
\end{verbatim}
whose last line reports the number of nodes, solutions and open branches.

\section{Changes from version 2}\label{sec:changes}

\begin{itemize}[leftmargin=*]
\item Proposition~1 of version~2 (every prime divisor of a unitary perfect
number is 3-Higgs) is withdrawn: its proof used the false divisibility
$p-1\mid p^e+1$.  The statement is now Hypothesis~H3.  The elimination of the
impostor kernels (version~2's Theorem~3 and Corollary~4) is stated under it
(Theorem~\ref{thm:conditional}), with a one-line argument for three of the five
kernels, and the source-kernel classification is no longer relied on.
\item Filter O and its monotonicity lemma are withdrawn
(Section~\ref{sec:higgs}).  The ten candidates it was used for are settled
unconditionally (Proposition~\ref{prop:ten}).
\item The exponent formula~\eqref{eq:exponent} is corrected: the seed term was
missing.
\item New: the exact search and Theorem~\ref{thm:intro}
(Section~\ref{sec:search}).
\item Section~\ref{sec:Heven}: the conjectural programme of version~2 is
removed.  Its hybrid semigroup-friable conjecture is false: if $k$ is an odd
prime in $\mathcal P_3$ and $r=4k+1$ is prime, then $\operatorname{ord}_r(2)=4k$
and $(r-1)/(4k)=1$, and such $k$ are plentiful.  Its log-mass conjecture is
false as stated, its estimate of the number of 3-Higgs primes up to $2^{44}$
was wrong, and its explicit form of conditional finiteness was not proved.
Smaller corrections: the largest witness has $115$ digits, not $85$; the
factor used for $m=1202$ has $87$ digits; the primality tests used are stated;
and the Cunningham table entry for $2,2426\mathrm L$ has the algebraic factor
$5$.  In the ``Status at $\max_a=10000$'' discussion of version~2, the
candidates $a=6759$ and $a=7338$ were described as closed by filter O; in the
final certificate run they are closed by filter N, which is what
Theorem~\ref{thm:conditional} uses.
\end{itemize}

\section*{Acknowledgments}

I thank the maintainers of FactorDB~\cite{FactorDB}, whose data underpin
filter N and Section~\ref{sec:Heven}.

\begin{thebibliography}{99}

\bibitem{SW1966}
M.~V.~Subbarao and L.~J.~Warren,
``Unitary perfect numbers,''
\emph{Canad. Math. Bull.} \textbf{9} (1966), 147--153.

\bibitem{Wall1975}
C.~R.~Wall,
``The fifth unitary perfect number,''
\emph{Canad. Math. Bull.} \textbf{18} (1975), 115--122.

\bibitem{Subbarao1970}
M.~V.~Subbarao,
``Are there an infinity of unitary perfect numbers?''
\emph{Amer. Math. Monthly} \textbf{77} (1970), 389--390.

\bibitem{Zsigmondy}
K.~Zsigmondy,
``Zur Theorie der Potenzreste,''
\emph{Monatsh. Math. Phys.} \textbf{3} (1892), 265--284.

\bibitem{SCNW1972}
M.~V.~Subbarao, T.~J.~Cook, R.~S.~Newberry, and J.~M.~Weber,
``On unitary perfect numbers,''
\emph{Delta} \textbf{3} (1972), no.~1, 22--26.

\bibitem{Wall1987}
C.~R.~Wall,
``On the largest odd component of a unitary perfect number,''
\emph{Fibonacci Quart.} \textbf{25} (1987), 312--316.

\bibitem{Wall1988}
C.~R.~Wall,
``New unitary perfect numbers have at least nine odd components,''
\emph{Fibonacci Quart.} \textbf{26} (1988), no.~4, 312--317.

\bibitem{Graham1989}
S.~W.~Graham,
``Unitary perfect numbers with squarefree odd part,''
\emph{Fibonacci Quart.} \textbf{27} (1989), 317--322.

\bibitem{Ford2014}
K.~Ford,
``Sieving very thin sets of primes, and Pratt trees with missing
primes,''
\emph{Int.~Math.~Res.~Not.~IMRN} \textbf{2014}, no.~11, 2955--2971.
(arXiv:1212.3498)

\bibitem{Wagstaff2002}
S.~S.~Wagstaff, Jr.,
\emph{The Cunningham Project},
\url{https://homes.cerias.purdue.edu/~ssw/cun/} (last accessed 2026).
See also Brillhart, Lehmer, Selfridge, Tuckerman, and Wagstaff,
\emph{Factorizations of $b^n \pm 1$, $b=2,3,5,6,7,10,11,12$ up to high
powers}, AMS, 3rd ed., 2002.

\bibitem{OEIS}
OEIS Foundation Inc.,
``The On-Line Encyclopedia of Integer Sequences,''
Sequence A002827 (unitary perfect numbers) and A057447 (3-Higgs primes),
\url{https://oeis.org/}.

\bibitem{FactorDB}
FactorDB online factor database, \url{https://factordb.com/}.

\bibitem{FKL}
K.~Ford, S.~V.~Konyagin, and F.~Luca,
``Prime chains and Pratt trees,''
\emph{Geom. Funct. Anal.} \textbf{20} (2010), 1231--1258;
arXiv:0904.0473.

\bibitem{BurrisYeats}
S.~Burris and K.~Yeats,
``The saga of the high school identities,''
\emph{Algebra Universalis} \textbf{52} (2005), 325--342.
Introduces the recursive sequence $\Sigma_a$ of exponent-$a$ Higgs
primes (motivational context for OEIS A057447); identifies $a > 1$
as an open infinitude/density problem.

\bibitem{WimmerNoschinski}
S.~Wimmer and L.~Noschinski,
``Pratt's primality certificates,''
\emph{Archive of Formal Proofs} (2013),
\url{https://www.isa-afp.org/entries/Pratt_Certificate.html}.

\end{thebibliography}

\end{document}

